\begin{claim}{C.62}{Proved}{Every equal-level prefix interval is contained in a compact three-word image}
Retain the original based plane loop $f_t$ and its generated
group $G$. On every open orbit under consideration choose one
actual global quotient, the original universal source cover
$p_O:X_O\to O$, and the specified real target lift
$F_O:X_O\to\R$ of C.43--C.44. Write
$\widetilde f_t^O$ for the based source lift.
For $x\in O$, choose $\widetilde x\in p_O^{-1}(x)$ and
$s,t\in[0,1]$ satisfying the exact real equality
\[
 F_O(\widetilde f_s^O\widetilde x)
   =F_O(\widetilde f_t^O\widetilde x).
                                                               \tag{C.62.1}
\]
Let $I$ be the compact interval in that proper real fiber
joining the two displayed points. Then
\[
 \begin{split}
 I&\subset
 \{\widetilde f_u^O(\widetilde f_v^O)^{-1}
                   \widetilde f_a^O\widetilde x:
                          0\leq a\leq v\leq u\leq1\},\\
 p_O(I)&\subset\mathcal W_3(x),\qquad
 \mathcal W_3(x)=
 \{f_u f_v^{-1}f_a(x):0\leq a\leq v\leq u\leq1\}.
 \end{split}                                                   \tag{C.62.2}
\]
The same assertion holds for the original inverse path,
with projected word $f_u^{-1}f_vf_a^{-1}$.

Consequently, let $x_n\in O_n$ tend to any $x\in\R^2$,
allow the open orbit and the chosen actual quotient to vary,
and let $J_n=p_{O_n}(I_n)$ be any intervals obtained from
(C.62.1). These intervals remain in one compact subset of
the plane and satisfy
\[
 \sup_{y\in J_n}\operatorname{dist}
                (y,\mathcal W_3(x))\longrightarrow0.
                                                               \tag{C.62.3}
\]
Every accumulation point therefore lies in the same
$G$-orbit as $x$. In particular, at a point fixed by $G$,
\[
                    \sup_{y\in J_n}|y-x|\longrightarrow0.
                                                               \tag{C.62.4}
\]
All source and target covering labels in these statements
are the original ones. In particular, (C.62.1) cannot be
replaced by equality modulo the deck group of a circle
target.
\end{claim}

\begin{proof}
\emph{The actual source chart and its two sides.}
C.43 provides a homeomorphism
\[
 \chi:X_O\longrightarrow\R^2,
       \qquad F_O=\operatorname{pr}_2\chi.
                                                               \tag{C.62.5}
\]
We use this exact chart for the proof and retain $p_O$.
Thus a real fiber is a complete horizontal line, and its
two complementary components are precisely its strict
sublevel and superlevel sets. Every original source lift
preserves the fiber partition. Its induced real target
map is increasing, because it is joined to the identity
by the based lifted path. The following argument uses
these actual homeomorphisms and their actual inverse
maps; it does not assume that the auxiliary product flow
from C.43 commutes with them.

Swap $s,t$ if necessary so that $s\leq t$, and put
\[
 \gamma(r)=\widetilde f_r^O\widetilde x,
       \qquad c=F_O\gamma(s)=F_O\gamma(t),
       \qquad
 Z=\{r\in[s,t]:F_O\gamma(r)=c\}.
                                                               \tag{C.62.6}
\]
The set $Z$ is compact and contains both endpoints.
Every complementary interval $(a,b)$ has endpoints
in $Z$, and the continuous function
$F_O\gamma(r)-c$ has one strict sign throughout it.

\emph{A single excursion is filled by an actual word triangle.}
Fix such an interval $(a,b)$. Suppress the superscript
$O$ and define
\[
 A=\widetilde f_a\widetilde x,\qquad
 h_r=\widetilde f_r\widetilde f_a^{-1}
       \quad(a\leq r\leq b),\qquad
 g=h_b,\qquad B=gA.
                                                               \tag{C.62.7}
\]
In particular $h_a=\id$, $A,B\in F_O^{-1}(c)$, and
$h_rA=\gamma(r)$.
Write $\beta_r$ for the increasing target map induced
by $h_r$. Suppose first that the excursion is positive.
Then
\[
 \beta_r(c)>c\quad(a<r<b),\qquad
 \beta_r^{-1}(c)<c,\qquad \beta_b(c)=c.
                                                               \tag{C.62.8}
\]
Since $\beta_b$ is increasing and fixes $c$, the two
paths
\[
       r\longmapsto h_rA,\qquad
       r\longmapsto g h_r^{-1}A
                                                               \tag{C.62.9}
\]
lie strictly on opposite sides of the fiber for
$a<r<b$. The first goes from $A$ to $B$ and the second
from $B$ to $A$. A negative excursion reverses their
two sides and changes nothing else in the argument.

For completeness, these side conditions give a precise
winding calculation even when either path has
self-intersections. Let $q$ be any point in the open
fiber interval between $A$ and $B$, with $A\ne B$.
In the chart (C.62.5), each path can be homotoped,
relative to its endpoints, to a two-segment path in
its own strict half-plane: use straight interpolation
with the polygonal path whose middle vertex is
directly above or below the midpoint of $A,B$,
at unit distance from the fiber on the specified
side. Throughout the interpolation
all interior parameter values stay strictly on that
side and hence avoid $q$; the endpoints also avoid
$q$. The resulting two polygonal paths form the
boundary of a diamond containing $q$. Their
concatenation therefore has winding $1$ or $-1$
about $q$. This proves the same nonzero winding for
(C.62.9), without imposing simplicity on the original
paths.

Now take the closed triangle
\[
 T_{a,b}=\{(u,v):a\leq v\leq u\leq b\}
\]
and the continuous map
\[
 \begin{split}
 \Psi(u,v)&=h_u h_v^{-1}A\\
          &=\widetilde f_u\widetilde f_v^{-1}
                                      \widetilde f_a\widetilde x.
 \end{split}                                                   \tag{C.62.10}
\]
On its edge $v=a$ this is the first path in (C.62.9);
on its edge $u=b$ it is the second path, in the order
specified there; and on the diagonal $u=v$ it is
the constant $A$. Thus its boundary has the nonzero
winding just computed about every such $q$.
If $q$ were missing from $\Psi(T_{a,b})$, this very
triangle would null-homotope that boundary in
$X_O\setminus\{q\}$, contradicting the winding.
The endpoints are already images of vertices.
When $A=B$ the interval is a singleton and the
conclusion is immediate. We have proved in all
cases that
\[
 [\gamma(a),\gamma(b)]_{F_O=c}
       \subset \Psi(T_{a,b}).
                                                               \tag{C.62.11}
\]
The equality in (C.62.10) follows by cancellation
of the displayed original lifts, so its three
letters have exactly the labels claimed in
(C.62.2).

\emph{The full zero set introduces no missing fiber intervals.}
Write $\chi\gamma(r)=(\xi(r),\zeta(r))$.
On $Z$ retain the real function $\xi|_Z$.
On every gap $(a,b)$ of $Z$, extend it by the
linear function joining $\xi(a)$ to $\xi(b)$;
denote the resulting function on $[s,t]$ by $e$.
This extension is continuous, including at
accumulation points of gaps. Indeed, on a gap of
length less than $\delta$, uniform continuity of
$\xi$ bounds the difference between $e(r)$ and
$\xi(r)$ by its modulus of continuity at
$\delta$. There are only finitely many gaps of
length at least $\delta$, and the linear extension
is continuous at both endpoints of each of
those gaps. Letting $\delta$ decrease proves the
assertion at every point of $Z$.

The continuous image $e([s,t])$ contains the
entire real interval between $\xi(s)$ and
$\xi(t)$. At a parameter in $Z$ its value gives
the original point $\gamma(r)$ on the fiber.
On a gap its value lies between the two endpoint
fiber coordinates, and so lies in the interval
already covered by (C.62.11). Consequently the
full interval $I$ is covered by the zero-level
track and these excursion intervals.
A zero-level track point is itself a three-letter
value, using $a=v=u=r$. All excursion triangles
use $0\leq a\leq v\leq u\leq1$.
This proves the first inclusion in (C.62.2).

\emph{Projection and the original deck equations.}
For each source lift we have
$p_O\widetilde f_r^O=f_rp_O$, and the same
equation holds for its inverse. Projecting
(C.62.10) therefore gives exactly
$f_u f_v^{-1}f_a(x)$, proving the second
inclusion in (C.62.2).
No sheet selection enters this comparison.
More explicitly, each original deck map $\delta$
commutes with every $\widetilde f_r^O$, because
the deck-valued commutator path starts at the
identity. C.44 supplies the retained target
equation
\[
                   F_O\delta=F_O+\eta_O(\delta).
                                                               \tag{C.62.12}
\]
Replacing $\widetilde x$ by $\delta\widetilde x$
therefore carries the entire interval and every
word value to their $\delta$-images, with the
same translated real level. Projection removes
that same deck map. Equality only modulo a
circle period would not give the single
horizontal fiber or the two sides used in
(C.62.8), which explains the exact hypothesis
(C.62.1).
Applying the already proved argument to the
actual continuous path $(\widetilde f_r^O)^{-1}$
gives the stated inverse-path inclusion and its
projected word $f_u^{-1}f_vf_a^{-1}$.

\emph{Uniform ambient control, across arbitrary orbits.}
Let
\[
 \Delta_3=\{(a,v,u)\in[0,1]^3:a\leq v\leq u\},\qquad
 W((a,v,u),y)=f_u f_v^{-1}f_a(y).
                                                               \tag{C.62.13}
\]
This is one jointly continuous global map,
independent of all orbit charts and source covers.
The set $E=\{x\}\cup\{x_n:n\geq1\}$ is compact.
Hence $W(\Delta_3\times E)$ is compact and
contains every $J_n$ by (C.62.2).
Continuity on this compact product also gives
\[
 \sup_{\lambda\in\Delta_3}
      |W(\lambda,x_n)-W(\lambda,x)|\longrightarrow0.
                                                               \tag{C.62.14}
\]
The distance in (C.62.3) is bounded by this
quantity. Every value $W(\lambda,x)$ belongs
to $Gx$, so (C.62.3) confines all accumulation
to that same orbit and excludes ambient escape.
At a $G$-fixed point, every such value is $x$;
this proves (C.62.4). The identical compactness
argument applies to the inverse-path words.

This strengthens C.61 to compact containment
near the specified nonopen orbit and proves
the full interval estimate at fixed frontier
points. At a nonfixed orbit, the compact set
$\mathcal W_3(x)$ may contain more than one
point. The statement does not yet select the
specified median or exponential boundary
value, nor construct a jointly continuous
lift of the moving quotient deformation.
Those two remaining requirements are retained
in C.12.
\end{proof}
